\section{Background and Related Work}
\label{sec:related-work}

We sit at the intersection of Roofline-guided GPU performance engineering, execution-free performance prediction, and LLM-based software engineering.
The relevant prior work is broad, but the distinctions that matter for this study are fairly crisp.

\paragraph{Roofline-guided GPU triage.}
The Roofline model gives a compact way to reason about whether a kernel is limited primarily by computation or by data movement \cite{samuelwilliamsRooflineInsightfulVisual2009}.
Its horizontal axis, arithmetic intensity, is not a full performance model; it is an interpretable ratio that helps developers decide which optimization direction is worth pursuing first.
That interpretability is why we use \rai as a triage target rather than a more opaque prediction objective.
Our contribution is therefore not another Roofline model, but an empirical study of when LLMs can estimate a Roofline-relevant signal from software artifacts before target execution.

\paragraph{Execution-free GPU performance prediction.}
Long before LLMs, GPU researchers studied analytical and learned predictors that estimate performance without executing the target kernel.
Their inputs include static program structure, microarchitectural assumptions, PTX/SASS features, and compiler-derived summaries \cite{hongAnalyticalModelGPU2009,baghsorkhiAdaptivePerformanceModeling2010,alavaniPredictingExecutionTime2018,guerreiroGPUStaticModeling2019a,braunSimpleModelPortable2021,emondsCUDAsapStaticallyDeterminedExecution2023,tranAccurateLatencyPredictor2025}.
Related analytical and learned predictors also study transfer across workloads or devices and learned representations such as GNNs or hybrid LLM/GNN predictors.
These works establish that execution-free prediction is valuable, but they generally target runtime, throughput, power, energy, or static execution statistics rather than the profiler-derived \bandwidthbound/\computebound regime we study.
They also usually rely on task-specific modeling pipelines, compiler analyses, or learned predictors rather than zero-training prompting.
Our question is narrower and more software-engineering-oriented: what can a developer tool obtain from existing code/build artifacts using an unmodified LLM, and where should that tool fall back to cheaper static features or runtime profiling?

This difference in objective also explains why we do not benchmark directly against these predictors, and instead compare against purpose-built static baselines (Section~\ref{sec:study-design}).
The analytical and learned models above predict a different quantity---execution time \cite{hongAnalyticalModelGPU2009,baghsorkhiAdaptivePerformanceModeling2010,alavaniPredictingExecutionTime2018}, latency \cite{tranAccurateLatencyPredictor2025}, power and energy under frequency scaling \cite{guerreiroGPUStaticModeling2019a,braunSimpleModelPortable2021}, or static basic-block frequencies \cite{emondsCUDAsapStaticallyDeterminedExecution2023}---so repurposing them to per-precision DRAM-level Roofline classification would require re-implementing them against a target they were never designed or validated for, yielding a comparison that reflects our re-implementation rather than the published method.
Several also assume inputs outside our setting (PTX feature pipelines, fitted hardware-independent features) or target older architectures with no artifact retargetable to Ampere/Hopper, and most are CUDA- or OpenCL-only and so cannot address the OpenMP-offload half of our corpus.
Where the prior work is a \emph{learned} predictor over compiler-IR features \cite{guerreiroGPUStaticModeling2019a,braunSimpleModelPortable2021,tranAccurateLatencyPredictor2025}, we instead instantiate that class faithfully for our target with the grouped-cross-validated random-forest and extra-trees baselines of Section~\ref{sec:study-design}, rather than mis-running an artifact built for a different objective.

\paragraph{LLMs for GPU optimization and code reasoning.}
Recent work has shown that LLMs can generate or optimize accelerator code, especially when compilation, execution, testing, or profiling is available inside the loop \cite{gongLanguageModelsCode2025,ouyangKernelBenchCanLLMs2025,chenCUDALLMLLMsCan2025,pengPerfCodeGenImprovingPerformance2025}.
Other work studies compiler-oriented representations and accelerator-code benchmark suites.
That literature is directly relevant to ICSE because it frames GPU tuning as an AI-assisted software-development activity.
Our study complements those systems rather than competing with them.
We do not optimize kernels directly; we study whether an LLM can provide an earlier triage signal when the profiler loop itself is unavailable or too costly.

\paragraph{LLMs for performance prediction.}
The closest neighboring work asks LLMs to predict GPU metrics directly, and the nearest two are both \emph{fine-tuned}.
LLMPerf fine-tunes models for OpenCL kernel performance estimation \cite{nguyen-nhatLLMPerfGPUPerformance2024}, and Omniwise fine-tunes a model to predict several GPU metrics, including arithmetic intensity, execution-free, reporting over 90\% of predictions within 10\% \cite{wangOmniwisePredictingGPU2025}.
Omniwise is the closest result to ours, and we confront it directly rather than competing with it: it answers a different question (what a \emph{trained} model can predict), it targets AMD MI250/MI300X with no released model retargetable to the NVIDIA GPUs and counters we study, and so it cannot be run on our corpus.
Our study is the complementary zero-training counterpart---if anything, Omniwise's strong fine-tuned numbers make the off-the-shelf gap we characterize the interesting question.
Prior work by Bolet \textit{et al.} studies coarse \bandwidthbound/\computebound classification of parallel code \cite{boletCanLargeLanguage2025}; the present work generalizes that precursor from a single coarse label to a per-precision, decomposed Roofline signal and to the source-versus-compilation evidence question.
Our contribution therefore differs in three ways.
First, we study the zero-training regime with off-the-shelf models rather than fine-tuned predictors.
Second, we target a decomposed signal by asking models to predict launch dimensions, per-precision FLOP counts, and DRAM bytes, from which both the regime label and \rai follow.
Third, we explicitly compare \sourceonly reasoning with prompts that paste target-kernel SASS instructions into the context, isolating the gap between source-visible semantics and compiler-realized behavior.
That distinction matters because compiler lowering, instruction selection, and floating-point implementations can introduce work that is not syntactically obvious in source code \cite{alfredCompilersPrinciplesTechniques2007,IEEEStandardFloatingPoint2019}.

\paragraph{Static reasoning with code LLMs.}
Our task also relates to a broader line of work showing that code LLMs can reason about latent program behavior, but only imperfectly \cite{liuCodeExecutionPretrained2023,lyuLargeLanguageModels2024,jainLiveCodeBenchHolisticContamination2024,xieCOREBenchmarkingLLMs2025,pydaShortcomingsCodeLLMs2025}.
Estimating \rai requires more than recognizing syntax: the model must propagate launch parameters, recover thread-level work, account for mixed precision, and reason about DRAM-visible traffic instead of source-level loads and stores.
That makes \rai prediction a useful probe of execution-relevant code reasoning in a form that is directly actionable for software engineering.

\paragraph{Gap addressed by this study.}
Taken together, prior work leaves an open gap.
We know that static models can predict GPU behavior without execution, that LLM optimizers benefit from profiler feedback, and that fine-tuned LLM predictors can estimate GPU metrics.
What remains underexplored is whether \emph{off-the-shelf} LLMs can provide a useful execution-free triage signal from ordinary software artifacts alone.
We address that gap by studying quantitative \rai prediction, separating source-visible from post-compilation evidence, and characterizing the regimes where off-the-shelf prompting is accurate, merely useful for triage, or unreliable.
